% Physical unit 010; accepted Session 02.
\GordanUnitStart{010}{printed page 5}
\GordanTopFolio{5}
\[
 \delta\int\left\{\sqrt{dx^2+dy^2+dz^2}+\lambda f\right\}=0
\]
vel si pro coordinatis rectangulis \gq{$x,y,z$} ellipticis $l,\psi$ utimur
\[
 \delta\int \sqrt{(1-e^2\cos l)\,dl^2+\cos^2l\,d\psi^2}=0.
\]

\begin{center}\S\ 4.\end{center}
Hic, si usitatas analistorum rationes adhibemus, ad extremum efficitur:
\[
 {\cos^2l\,{d\psi\over dl}\over\sqrt{1-e^2\cos^2l+\cos^2l\,{d\psi\over dl}{d\psi\over dl}}}=\mathrm{Const}
\]
qua in formula adhibito:
\[
 {\cos l\,{d\psi\over dl}\over\sqrt{1-e^2\cos^2l+\cos^2l\,{d\psi\over dl}{d\psi\over dl}}}=\sin T
\]
\hfill conf. Cap. I. \S\ 9.

est
\[
 \cos l\sin T=\mathrm{Const}
\]
et si
\[
 T={\pi\over 2}\quad\hbox{sumto}\qquad l=l_0\quad\hbox{ponimus,}
\]
\[
 \cos l\sin T=\cos l_0,
\]
\[
 \cos l\cos T=\sqrt{\sin^2l_0-\sin^2l},
\qquad
 \tg T={\cos l_0\over\sqrt{\sin^2l_0-\sin^2l}}.
\]

\begin{center}\S\ 5.\end{center}
Hoc est; linea geodetica circulum aliquem tangens; tangit alium quoque circulum, qui situm illi convenientem habet. Porro, quod semper:
\[
 l_0>l
\]
est, nunquam fieri potest, ut linea geodetica ultra hos circulos ab aequatore recedat. Propterea lineam geodeticam circulum $l_0$ tangentem nominamus lineam geodeticam $l_0$ et punctum illud, quo circulus tangitur, exordium lineae geodeticae, unde deinceps punctorum et extensiones \gq{$s$} et longitudines $\psi$ metiemur.

\begin{center}\S\ 6.\end{center}
Ad quas e latitudine $l$ determinandas, formulis II ad formulas Cap. I. \S\ 9 propositas relatis, est:
